% GENERATED FILE. Edit canonical lessons, then run npm run generate:books.
% canonical-source-hash: fnv1a64:42c1308524057f10
% canonical-lessons: KA-C206-edavu, KA-C206-tevalu, KA-C206-himbalisu, KA-C206-tappisikollu, KA-C206-aturapadu, KA-R206-first-pass-words-from-chapters-198-201, KA-R206-second-pass-words-from-chapters-202-206

\chapter{To trip, To crawl, To follow, To escape, To hurry}
\label{ch:ka-to-trip-to-crawl-to-follow-to-escape-to-hurry}
\hlchaptermodality{206}

\begin{chapteropening}
\textbf{By the end of this chapter:} \emph{I can say to trip, to crawl, to follow, to escape and to hurry.}

Complete the last lesson of chapter 206: I can say to trip, to crawl, to follow, to escape and to hurry.
\end{chapteropening}

\section[\texorpdfstring{\kn{ಎಡವು} (eḍavu)}{ಎಡವು (eḍavu)}]{\kn{ಎಡವು} (eḍavu) — to trip, to stumble}

\label{lesson:KA-C206-edavu}



\begin{quote}
Before the new one: say the Kannada for to knead, then the Kannada for to roll out.
\end{quote}

\subsection*{You'll want to know: \kn{ಎಡವು}}
\textbf{\kn{ಎಡವು}} — \emph{eḍavu} — \textquotedblleft{}to trip, to stumble\textquotedblright{}.

\begin{grammarlens}[title={What you've built}]
One more word for this chapter.
\end{grammarlens}

\subsection*{Your turn}
\begin{itemize}
\raggedright
  \item \emph{Say it:} \emph{eḍavu}
  \item \emph{Say it:} \emph{eḍavu}, once more
  \item \emph{Say it:} say \emph{eḍavu}
  \item \emph{Recall:} say the Kannada for to knead, then the Kannada for to roll out, then say \emph{eḍavu} again
\end{itemize}

\begin{tcolorbox}[breakable,colback=teal!4,colframe=teal!35!black,title={Before you move on}]
What does \kn{ಎಡವು} mean? (\textquotedblleft{}To trip, to stumble\textquotedblright{}.) Say it once more.
\end{tcolorbox}
\section[\texorpdfstring{\kn{ತೆವಳು} (tevaḷu)}{ತೆವಳು (tevaḷu)}]{\kn{ತೆವಳು} (tevaḷu) — to crawl}

\label{lesson:KA-C206-tevalu}



\begin{quote}
Before the new one: say the Kannada for to roll out, then the Kannada for to trip.
\end{quote}

\subsection*{You'll want to know: \kn{ತೆವಳು}}
\textbf{\kn{ತೆವಳು}} — \emph{tevaḷu} — \textquotedblleft{}to crawl\textquotedblright{}.

\begin{grammarlens}[title={What you've built}]
One more word for this chapter.
\end{grammarlens}

\subsection*{Your turn}
\begin{itemize}
\raggedright
  \item \emph{Say it:} \emph{tevaḷu}
  \item \emph{Say it:} \emph{tevaḷu}, once more
  \item \emph{Say it:} say \emph{tevaḷu}
  \item \emph{Recall:} say the Kannada for to roll out, then the Kannada for to trip, then say \emph{tevaḷu} again
\end{itemize}

\begin{tcolorbox}[breakable,colback=teal!4,colframe=teal!35!black,title={Before you move on}]
What does \kn{ತೆವಳು} mean? (\textquotedblleft{}To crawl\textquotedblright{}.) Say it once more.
\end{tcolorbox}
\section[\texorpdfstring{\kn{ಹಿಂಬಾಲಿಸು} (himbālisu)}{ಹಿಂಬಾಲಿಸು (himbālisu)}]{\kn{ಹಿಂಬಾಲಿಸು} (himbālisu) — to follow (someone)}

\label{lesson:KA-C206-himbalisu}



\begin{quote}
Before the new one: say the Kannada for to trip, then the Kannada for to crawl.
\end{quote}

\subsection*{You'll want to know: \kn{ಹಿಂಬಾಲಿಸು}}
\textbf{\kn{ಹಿಂಬಾಲಿಸು}} — \emph{himbālisu} — \textquotedblleft{}to follow (someone)\textquotedblright{}.

\begin{grammarlens}[title={What you've built}]
One more word for this chapter.
\end{grammarlens}

\subsection*{Your turn}
\begin{itemize}
\raggedright
  \item \emph{Say it:} \emph{himbālisu}
  \item \emph{Say it:} \emph{himbālisu}, once more
  \item \emph{Say it:} say \emph{himbālisu}
  \item \emph{Recall:} say the Kannada for to trip, then the Kannada for to crawl, then say \emph{himbālisu} again
\end{itemize}

\begin{tcolorbox}[breakable,colback=teal!4,colframe=teal!35!black,title={Before you move on}]
What does \kn{ಹಿಂಬಾಲಿಸು} mean? (\textquotedblleft{}To follow (someone)\textquotedblright{}.) Say it once more.
\end{tcolorbox}
\section[\texorpdfstring{\kn{ತಪ್ಪಿಸಿಕೊಳ್ಳು} (tappisikoḷḷu)}{ತಪ್ಪಿಸಿಕೊಳ್ಳು (tappisikoḷḷu)}]{\kn{ತಪ್ಪಿಸಿಕೊಳ್ಳು} (tappisikoḷḷu) — to escape, to dodge}

\label{lesson:KA-C206-tappisikollu}



\begin{quote}
Before the new one: say the Kannada for to crawl, then the Kannada for to follow.
\end{quote}

\subsection*{You'll want to know: \kn{ತಪ್ಪಿಸಿಕೊಳ್ಳು}}
\textbf{\kn{ತಪ್ಪಿಸಿಕೊಳ್ಳು}} — \emph{tappisikoḷḷu} — \textquotedblleft{}to escape, to dodge\textquotedblright{}.

\begin{grammarlens}[title={What you've built}]
One more word for this chapter.
\end{grammarlens}

\subsection*{Your turn}
\begin{itemize}
\raggedright
  \item \emph{Say it:} \emph{tappisikoḷḷu}
  \item \emph{Say it:} \emph{tappisikoḷḷu}, once more
  \item \emph{Say it:} say \emph{tappisikoḷḷu}
  \item \emph{Recall:} say the Kannada for to crawl, then the Kannada for to follow, then say \emph{tappisikoḷḷu} again
\end{itemize}

\begin{tcolorbox}[breakable,colback=teal!4,colframe=teal!35!black,title={Before you move on}]
What does \kn{ತಪ್ಪಿಸಿಕೊಳ್ಳು} mean? (\textquotedblleft{}To escape, to dodge\textquotedblright{}.) Say it once more.
\end{tcolorbox}
\section[\texorpdfstring{\kn{ಆತುರಪಡು} (āturapaḍu)}{ಆತುರಪಡು (āturapaḍu)}]{\kn{ಆತುರಪಡು} (āturapaḍu) — to hurry, to be in a rush}

\label{lesson:KA-C206-aturapadu}



\begin{quote}
Before the new one: say the Kannada for to follow, then the Kannada for to escape.
\end{quote}

\subsection*{You'll want to know: \kn{ಆತುರಪಡು}}
\textbf{\kn{ಆತುರಪಡು}} — \emph{āturapaḍu} — \textquotedblleft{}to hurry, to be in a rush\textquotedblright{}.

\begin{grammarlens}[title={What you've built}]
One more word for this chapter.
\end{grammarlens}

\subsection*{Your turn}
\begin{itemize}
\raggedright
  \item \emph{Say it:} \emph{āturapaḍu}
  \item \emph{Say it:} \emph{āturapaḍu}, once more
  \item \emph{Say it:} say \emph{āturapaḍu}
  \item \emph{Recall:} say the Kannada for to follow, then the Kannada for to escape, then say \emph{āturapaḍu} again
\end{itemize}

\begin{tcolorbox}[breakable,colback=teal!4,colframe=teal!35!black,title={Before you move on}]
What does \kn{ಆತುರಪಡು} mean? (\textquotedblleft{}To hurry, to be in a rush\textquotedblright{}.) Say it once more.
\end{tcolorbox}
\section[(dialogue)]{Words from chapters 198-201 — a first pass}

\label{lesson:KA-R206-first-pass-words-from-chapters-198-201}



\begin{quote}
Say the words for to escape and to hurry.
\end{quote}

\subsection*{Your turn}
Say these aloud:

\begin{itemize}
\raggedright
  \item to scrub, to bring, to carry away, to hold, to let go
  \item to swallow, to spread out, to fold, to iron, to change
  \item to get well, to itch, to die, to finish, to start
  \item to postpone, to explain, to pass, to remember, to print
\end{itemize}

\begin{tcolorbox}[breakable,colback=teal!4,colframe=teal!35!black,title={Before you move on}]
To scrub? (\textbf{\kn{ಉಜ್ಜು}}, \emph{ujju}.) To bring? (\textbf{\kn{ತರು}}, \emph{taru}.) To carry away? (\textbf{\kn{ಒಯ್ಯು}}, \emph{oyyu}.) To hold? (\textbf{\kn{ಹಿಡಿ}}, \emph{hiḍi}.) To let go? (\textbf{\kn{ಬಿಡು}}, \emph{biḍu}.) To swallow? (\textbf{\kn{ನುಂಗು}}, \emph{nuṅgu}.) To spread out? (\textbf{\kn{ಹಾಸು}}, \emph{hāsu}.) To fold? (\textbf{\kn{ಮಡಚು}}, \emph{maḍacu}.) To iron? (\textbf{\kn{ಇಸ್ತ್ರಿ} \kn{ಮಾಡು}}, \emph{istri māḍu}.) To change? (\textbf{\kn{ಬದಲಿಸು}}, \emph{badalisu}.) To get well? (\textbf{\kn{ಗುಣವಾಗು}}, \emph{guṇavāgu}.) To itch? (\textbf{\kn{ತುರಿಸು}}, \emph{turisu}.) To die? (\textbf{\kn{ಸಾಯು}}, \emph{sāyu}.) To finish? (\textbf{\kn{ಮುಗಿಸು}}, \emph{mugisu}.) To start? (\textbf{\kn{ಶುರು} \kn{ಮಾಡು}}, \emph{śuru māḍu}.) To postpone? (\textbf{\kn{ಮುಂದೂಡು}}, \emph{mundūḍu}.) To explain? (\textbf{\kn{ವಿವರಿಸು}}, \emph{vivarisu}.) To pass? (\textbf{\kn{ಪಾಸಾಗು}}, \emph{pāsāgu}.) To remember? (\textbf{\kn{ನೆನಪಿಡು}}, \emph{nenapiḍu}.) To print? (\textbf{\kn{ಮುದ್ರಿಸು}}, \emph{mudrisu}.)
\end{tcolorbox}
\section[(dialogue)]{Words from chapters 202-206 — a second pass}

\label{lesson:KA-R206-second-pass-words-from-chapters-202-206}



\begin{quote}
Say the words for to complain and to register.
\end{quote}

\subsection*{Your turn}
Say these aloud:

\begin{itemize}
\raggedright
  \item to complain, to register, to continue, to contact, to discuss
  \item to grow, to bloom, to wilt, to blow, to melt
  \item to publish, to arrange, to measure, to have fun, to compete
  \item to pluck, to grate, to strain, to knead, to roll out
  \item to trip, to crawl, to follow, to escape, to hurry
\end{itemize}

\begin{tcolorbox}[breakable,colback=teal!4,colframe=teal!35!black,title={Before you move on}]
To complain? (\textbf{\kn{ದೂರು} \kn{ಕೊಡು}}, \emph{dūru koḍu}.) To register? (\textbf{\kn{ನೋಂದಾಯಿಸು}}, \emph{nōndāyisu}.) To continue? (\textbf{\kn{ಮುಂದುವರಿಸು}}, \emph{munduvarisu}.) To contact? (\textbf{\kn{ಸಂಪರ್ಕಿಸು}}, \emph{samparkisu}.) To discuss? (\textbf{\kn{ಚರ್ಚಿಸು}}, \emph{carcisu}.) To grow? (\textbf{\kn{ಬೆಳೆಸು}}, \emph{beḷesu}.) To bloom? (\textbf{\kn{ಅರಳು}}, \emph{araḷu}.) To wilt? (\textbf{\kn{ಬಾಡು}}, \emph{bāḍu}.) To blow? (\textbf{\kn{ಬೀಸು}}, \emph{bīsu}.) To melt? (\textbf{\kn{ಕರಗು}}, \emph{karagu}.) To publish? (\textbf{\kn{ಪ್ರಕಟಿಸು}}, \emph{prakaṭisu}.) To arrange? (\textbf{\kn{ಜೋಡಿಸು}}, \emph{jōḍisu}.) To measure? (\textbf{\kn{ಅಳೆ}}, \emph{aḷe}.) To have fun? (\textbf{\kn{ಮಜಾ} \kn{ಮಾಡು}}, \emph{majā māḍu}.) To compete? (\textbf{\kn{ಸ್ಪರ್ಧಿಸು}}, \emph{spardhisu}.) To pluck? (\textbf{\kn{ಕೀಳು}}, \emph{kīḷu}.) To grate? (\textbf{\kn{ತುರಿ}}, \emph{turi}.) To strain? (\textbf{\kn{ಸೋಸು}}, \emph{sōsu}.) To knead? (\textbf{\kn{ನಾದು}}, \emph{nādu}.) To roll out? (\textbf{\kn{ಲಟ್ಟಿಸು}}, \emph{laṭṭisu}.) To trip? (\textbf{\kn{ಎಡವು}}, \emph{eḍavu}.) To crawl? (\textbf{\kn{ತೆವಳು}}, \emph{tevaḷu}.) To follow? (\textbf{\kn{ಹಿಂಬಾಲಿಸು}}, \emph{himbālisu}.) To escape? (\textbf{\kn{ತಪ್ಪಿಸಿಕೊಳ್ಳು}}, \emph{tappisikoḷḷu}.) To hurry? (\textbf{\kn{ಆತುರಪಡು}}, \emph{āturapaḍu}.)
\end{tcolorbox}
